% Part: incompleteness
% Chapter: representability-in-q
% Section: sigma1-completeness

\documentclass[../../../include/open-logic-section]{subfiles}

\begin{document}

\olfileid{inc}{inp}{s1c}
\olsection{\texorpdfstring{$\Sigma_1$}{सिग्मा-१} संपूर्णता}

$\Th{Q}$ आणि तिच्या सुसंगत, स्वयंसिद्धकीकरण करता येणाऱ्या
विस्तारांची अपूर्णता असूनही, संख्यांकांविषयी अनेक मूलभूत तथ्ये
$\Th{Q}$ सिद्ध करते, हे आपण पाहिले आहे. ही क्षमता प्रत्यक्षात
आणखी बरीच व्यापक आहे. $\Th{Q}$ मध्ये काय सिद्ध करता येते
याची व्याप्ती समजण्यासाठी $\Delta_0$, $\Sigma_1$ आणि
$\Pi_1$ !!{formula}s यांच्या संकल्पना मांडू. स्थूलमानाने,
$\Sigma_1$ !!{formula} हे $\lexists[x][!B(x)]$ या रूपाचे
असते; येथे $!B$ फक्त विधानीय संयोजके आणि परिबद्ध संख्यापक
वापरून तयार केलेले असते. $!A$ हे $\Struct{N}$ मध्ये सत्य
असलेले $\Sigma_1$ !!{sentence} असेल, तर
$\Th{Q} \Proves !A$, हे आपण दाखवू
(\olref{thm:sigma1-completeness}).

\begin{defn}
\ollabel{defn:bd-quant}
\emph{परिबद्ध अस्तित्ववाची !!{formula}} हे
$\lexists[x][(x < t \land !A(x))]$ या रूपाचे असते; येथे $t$
हे संबंधित संख्यापकाचे चल नसलेले कोणतेही पद आहे. परंपरेने ते
$\bexists{x < t}{!A(x)}$ असे लिहितो.
%
\emph{परिबद्ध वैश्विक !!{formula}} हे
$\lforall[x][(x < t \lif !A(x))]$ या रूपाचे असते; येथेही
$t$ या पदात संबंधित संख्यापकाचे चल नसते. परंपरेने ते
$\bforall{x < t}{!A(x)}$ असे लिहितो.
\end{defn}

\begin{defn}
\ollabel{defn:delta0-sigma1-pi1-frm}
!!{formula} $!B$ हे $\Delta_0$ असते, जर ते आण्विक
!!{formula}s पासून केवळ विधानीय संयोजके आणि परिबद्ध
संख्यकीकरण वापरून तयार झाले असेल.
%
!!{formula} $!A$ हे $\Sigma_1$ असते, जर
$!A \ident \lexists[x][!B(x)]$ आणि $!B$ हे $\Delta_0$
असेल.
%
!!{formula} $!A$ हे $\Pi_1$ असते, जर
$!A \ident \lforall[x][!B(x)]$ आणि $!B$ हे $\Delta_0$
असेल.
\end{defn}

\begin{lem}
\ollabel{lem:q-proves-clterm-id} समजा $t$ हे बंद पद असून
$\Value{t}{N} = n$. तर $\Th{Q} \Proves \eq[t][\num n]$.
\end{lem}

\begin{proof}
$t$ च्या गुंतागुंतीवर विगमन करू. आधार प्रकरणात
$\Value{\Obj 0}{N} = 0$, आणि
$\num 0 \ident \Obj 0$ असल्यामुळे
$\Th{Q} \Proves \eq[\Obj 0][\num 0]$.
%
विगमनाच्या टप्प्यासाठी $t_1$ आणि $t_2$ ही पदे अशी घ्या की
$\Value{t_1}{N} = n_1$, $\Value{t_2}{N} = n_2$,
$\Th{Q} \Proves \eq[t_1][\num n_1]$, आणि
$\Th{Q} \Proves \eq[t_2][\num n_2]$.

मग $\Value{(t_1')}{N} = n_1 + 1$. विगमन गृहीतक आणि
$\eq[\num{n_1}'][\num{n_1}']$ या !!{formula} वर
प्रथम-क्रम एकरूपतेचे नियम लावल्यास
$\Th{Q} \Proves \eq[t_1'][{\num n_1}']$ मिळते.
संख्यांकांच्या व्याख्येने
$\Th{Q} \Proves \eq[t_1'][\num{n_1 + 1}]$ मिळते.

बेरीजेसाठी,
\[
      \Value{(t_1 + t_2)}{N}
    = \Value{t_1}{N} + \Value{t_2}{N}
    = n_1 + n_2.
\]
विगमन गृहीतक आणि एकरूपतेच्या नियमांनी प्रथम
$\Th{Q} \Proves \eq[t_1 + t_2][\num{n_1} + t_2]$,
आणि ते नियम पुन्हा लावल्यास
$\Th{Q} \Proves \eq[t_1 + t_2][\num{n_1} + \num{n_2}]$
मिळते. \olref[inc][req][bre]{lem:q-proves-add} नुसार
$\Th{Q} \Proves \eq[\num{n_1} + \num{n_2}][\num{n_1 + n_2}]$;
म्हणून $\Th{Q} \Proves \eq[t_1 + t_2][\num{n_1 + n_2}]$.

$\times$ साठीही
\olref[inc][req][bre]{lem:q-proves-mult} वापरून असाच
युक्तिवाद करता येतो.
%
अंकगणितातील बंद पदांचे हे सर्व प्रकार आहेत. म्हणून
$\Value{t}{N} = n$ असलेल्या प्रत्येक बंद पद~$t$ साठी
$\Th{Q} \Proves \eq[t][\num n]$.
\end{proof}

\begin{prob}
\olref{lem:q-proves-clterm-id} मधील $\times$ संबंधीचा
भाग सविस्तर सिद्ध करा.
\end{prob}

\begin{lem}
\ollabel{lem:atomic-completeness}
समजा $t_1$ आणि $t_2$ ही बंद पदे आहेत. तर
\begin{enumerate}
\item $\Value{t_1}{N} = \Value{t_2}{N}$ असेल, तर
    $\Th{Q} \Proves \eq[t_1][t_2]$.
\item $\Value{t_1}{N} \neq \Value{t_2}{N}$ असेल, तर
    $\Th{Q} \Proves \eq/[t_1][t_2]$.
\item $\Value{t_1}{N} < \Value{t_2}{N}$ असेल, तर
    $\Th{Q} \Proves t_1 < t_2$.
\item $\Value{t_2}{N} \leq \Value{t_1}{N}$ असेल, तर
    $\Th{Q} \Proves \lnot(t_1 < t_2)$.
\end{enumerate}
\end{lem}

\begin{proof}
दिलेले $t_1$ आणि $t_2$ घेऊन
$n = \Value{t_1}{N}$ आणि $m = \Value{t_2}{N}$
निश्चित करू.

समजा $!A \ident t_1 = t_2$.
\olref{lem:q-proves-clterm-id} नुसार
$\Th{Q} \Proves \eq[t_1][\num n]$ आणि
$\Th{Q} \Proves \eq[t_2][\num m]$.
$n = m$ असेल, तर
$\Th{Q} \Proves \eq[\num n][\num m]$;
म्हणून एकरूपतेच्या संक्रमणशीलतेने
$\Th{Q} \Proves \eq[t_1][t_2]$.
$n \neq m$ असेल, तर
$\Th{Q} \Proves \eq/[\num n][\num m]$;
एकरूपतेचे नियम पुन्हा वापरल्यास
$\Th{Q} \Proves \eq/[t_1][t_2]$.

आता $!A \ident t_1 < t_2$ असे घेऊ. दोन्ही प्रकरणांत
$!Q_8$ या स्वयंसिद्धकाचा आधार घेऊ. त्यानुसार सर्व
$x,y$ साठी
$x < y \liff \lexists[z][\eq[z' + x][y]]$.

समजा $\Sat{N}{t_1 < t_2}$. मग असा $k \in \Nat$
असतो की $n + k + 1 = m$.
\olref{lem:q-proves-clterm-id} नुसार
$\Th{Q} \Proves \eq[t_1][\num n]$ आणि
$\Th{Q} \Proves \eq[t_2][\num m]$.
या प्रमेयाच्या पहिल्या भागाने निश्चित संख्यांकांची बेरीज
काढल्यास
$\Th{Q} \Proves \eq[{\num k}' + \num n][\num m]$.
एकरूपतेच्या संक्रमणशीलतेने मग
$\Th{Q} \Proves \eq[{\num k}' + t_1][t_2]$;
म्हणून
$\Th{Q} \Proves \lexists[z][\eq[z' + t_1][t_2]]$.
$!Q_8$ ची उजवीकडून डावीकडे दिशा वापरल्यास
$\Th{Q} \Proves t_1 < t_2$.

याउलट $\Sat/{N}{t_1 < t_2}$, म्हणजे $m \leq n$,
असे समजा.
%
$\Th{Q}$ मध्ये काम करून $t_1 < t_2$ असे गृहीत धरा.
$!Q_8$ ची डावीकडून उजवीकडे दिशा लावल्यास असा $z$
मिळतो की $\eq[z' + t_1][t_2]$.
$\Th{Q} \Proves \eq[t_1][\num n]$ आणि
$\Th{Q} \Proves \eq[t_2][\num m]$ असल्यामुळे
$\eq[z' + \num n][\num m]$.
%
$m$ वर बाह्य विगमन करू: प्रत्येक पायरीवर $!Q_5$
वापरून बेरीज उत्तरवर्तीच्या रूपात लिहू आणि उत्तरवर्तीचे
एकास-एक असणे वापरून दोन्ही बाजूंचे एक-एक उत्तरवर्ती
रद्द करू. अशा रीतीने
$\eq[z' + \num{n - m}][\Obj 0]$ मिळते.
$m = n$ असेल, तर शून्य जोडण्याच्या स्वयंसिद्धकाने
$\eq[z'][\Obj 0]$ मिळते; पण $!Q_2$ नुसार उत्तरवर्ती
शून्य नसतो, म्हणून विरोधाभास मिळतो.
$m < n$ असेल, तर $!Q_5$ पुन्हा वापरून
$\eq[(z' + \num{n - m - 1})'][\Obj 0]$ मिळते;
$!Q_2$ मुळे पुन्हा विरोधाभास मिळतो.
म्हणून $\Th{Q} \Proves \lnot(t_1 < t_2)$.
\end{proof}

\begin{lem}
\ollabel{lem:bounded-quant-equiv}
समजा $!A$ हे !!a{formula}, $t$ हे बंद पद, आणि
$k=\Value{t}{N}$. तर
\begin{enumerate}
\item $\Th{Q} \Proves \bforall{x<t}{!A(x)}$ तेव्हाच आणि तरच $\Th{Q} \Proves
    !A(\num 0) \land \dots \land !A(\num{k-1})$.
\item $\Th{Q} \Proves \bexists{x<t}{!A(x)}$ तेव्हाच आणि तरच $\Th{Q} \Proves
    !A(\num 0) \lor \dots \lor !A(\num{k-1})$.
\end{enumerate}
\end{lem}

\begin{proof}
    परिबद्ध वैश्विक संख्यापकाचे प्रकरण सिद्ध करू.
    $\Value{t}{N} = 0$ असेल, तर
    \olref[inc][req][min]{lem:less-zero} नुसार
    $x<\num 0$ अशी कोणतीही संख्या नसते; म्हणून
    समतुल्यतेची डावी बाजू $\Th{Q}$ मध्ये सिद्ध होते.
    उजवीकडील रिक्त संयोगाला $\ltrue$ मानतो;
    तेही $\Th{Q}$ मध्ये सिद्ध होते.
    %
    म्हणून आता एखाद्या नैसर्गिक संख्या~$k$ साठी
    $\Value{t}{N} = k+1$ असे समजू.
    \olref{lem:q-proves-clterm-id} नुसार
    $\bforall{x<\num{k+1}}{!A(x)}$ या रूपाच्या
    !!a{formula} साठी काम करणे पुरेसे आहे.

    समजा
    $\Th{Q} \Proves \bforall{x<\num{k+1}}{!A(x)}$,
    आणि $n \leq k$ घेऊ.
    \olref{lem:atomic-completeness} नुसार
    $\Th{Q} \Proves \num n < \num{k+1}$;
    म्हणून तर्कनियमांनी
    $\Th{Q} \Proves !A(\num n)$.
    प्रत्येक $n \leq k$ साठी हे $k+1$ वेळा
    लावल्यास अपेक्षेप्रमाणे
    $\Th{Q} \Proves !A(\num 0)
    \land \dots \land !A(\num k)$ मिळते.

    उलट दिशेसाठी समजा
    $\Th{Q} \Proves
    !A(\num 0) \land \dots \land !A(\num k)$.
    $\Th{Q}$ मध्ये $x < \num{k+1}$ असे समजू.
    \olref[inc][req][min]{lem:less-nsucc} नुसार
    $x = \num 0 \lor \dots \lor x = \num k$.
    मग तर्कनियमांनी $!A(x)$ मिळते आणि त्यावरून
    $\bforall{x<\num{k+1}}{!A(x)}$ हे वैश्विक विधान
    मिळते.

    परिबद्ध अस्तित्ववाची संख्यकीकरण असलेल्या
    !!{formula}s साठी समतुल्यतेचा पुरावाही असाच आहे.
\end{proof}

\begin{prob}
\olref{lem:bounded-quant-equiv} मधील अस्तित्ववाची
प्रकरणाचा सविस्तर पुरावा द्या.
\end{prob}

\begin{lem}
\ollabel{lem:delta0-completeness}
$!A$ हे $\Struct{N}$ मध्ये सत्य असलेले
$\Delta_0$ !!{sentence} असेल, तर
$\Th{Q} \Proves !A$.
\end{lem}

\begin{proof}
!!{formula} च्या गुंतागुंतीवर विगमन करू.
%
आधार प्रकरण \olref{lem:atomic-completeness} ने
मिळते; आता विगमनाचा टप्पा पाहू. सोयीसाठी,
निषेधाच्या प्रकरणाची विभागणी निषिद्ध !!{formula}
च्या रचनेनुसार करू. इतर विधानीय संयोजक पुढील
प्रकरणांमधून संक्षेपाने व्यक्त करता येतात.

\begin{enumerate}
\item $(!A \land !B)$ हे $\Struct{N}$ मध्ये सत्य
असेल, तर $!A$ आणि $!B$ दोन्ही $\Struct{N}$ मध्ये सत्य आहेत.
विगमन गृहीतकाने
$\Th{Q} \Proves !A$ आणि
$\Th{Q} \Proves !B$; म्हणून तर्कनियमांनी
$\Th{Q} \Proves (!A \land !B)$.
%
\item $\lnot (!A \land !B)$ हे $\Struct{N}$
मध्ये सत्य असेल, तर $\lnot !A$ किंवा
$\lnot !B$ हे $\Struct{N}$ मध्ये सत्य आहे. सामान्यत्वाला बाधा
न आणता पहिले सत्य आहे असे धरू.
विगमन गृहीतकाने $\Th{Q} \Proves \lnot !A$;
म्हणून तर्कनियमांनी
$\Th{Q} \Proves \lnot (!A \land !B)$.
%
\item $(!A \lor !B)$ हे $\Struct{N}$ मध्ये सत्य
असेल, तर $!A$ हे $\Struct{N}$ मध्ये सत्य आहे किंवा
$!B$ हे $\Struct{N}$ मध्ये सत्य आहे.
सामान्यत्वाला बाधा न आणता पहिले सत्य आहे
असे धरू. विगमन गृहीतकाने
$\Th{Q} \Proves !A$;
म्हणून तर्कनियमांनी
$\Th{Q} \Proves (!A \lor !B)$.
%
\item $\lnot(!A \lor !B)$ हे $\Struct{N}$
मध्ये सत्य असेल, तर $\lnot !A$ आणि
$\lnot !B$ दोन्ही $\Struct{N}$ मध्ये सत्य आहेत.
विगमन गृहीतकाने
$\Th{Q} \Proves \lnot !A$ आणि
$\Th{Q} \Proves \lnot !B$.
त्यामुळे तर्कनियमांनी
$\Th{Q} \Proves \lnot(!A \lor !B)$.
%
\item $\bforall{x<t}{!A(x)}$ हे
$\Struct{N}$ मध्ये सत्य आहे असे समजा;
येथे $t$ हे बंद पद आणि $k=\Value{t}{N}$.
सर्व $n < \Value{t}{N}$ साठी
$!A(\num n)$ हे $\Struct{N}$ मध्ये सत्य
असेल, तर विगमन गृहीतक आणि तर्कनियमांनी
$\Th{Q} \Proves
!A(\num 0) \land \dots \land !A(\num{k-1})$.
\olref{lem:bounded-quant-equiv} नुसार
$\Th{Q} \Proves \bforall{x<t}{!A(x)}$.
%
\item परिबद्ध अस्तित्ववाची संख्यापकाचे
प्रकरणही वैश्विक संख्यापकासारखेच आहे:
येथे $\bexists{x < t}{!A(x)}$ या रूपाचे
!!a{sentence} असते.
%
\item $\lnot \bforall{x<t}{!A(x)}$ हे
$\Struct{N}$ मध्ये सत्य आहे असे समजा;
येथे $t$ हे बंद पद. हे !!{sentence}
$\bexists{x<t}{\lnot !A(x)}$ या !!{sentence}
शी समतुल्य आहे आणि ही समतुल्यता
$\Th{Q}$ मध्येही सिद्ध करता येते. म्हणून
परिबद्ध अस्तित्ववाची संख्यापकाचा युक्तिवाद
लागू करता येतो.
%
\item त्याचप्रमाणे, $t$ हे बंद पद असताना
$\lnot \bexists{x<t}{!A(x)}$ हे
$\Struct{N}$ मध्ये सत्य आहे असे समजा.
हे !!{sentence} $\Th{Q}$ मध्ये
$\bforall{x<t}{\lnot!A(x)}$ शी समतुल्य
आहे; म्हणून परिबद्ध वैश्विक संख्यापकाचा
युक्तिवाद लागू करता येतो.
%
\item अखेरीस $\lnot !A$ हे
$\Struct{N}$ मध्ये सत्य आहे असे समजा.
उरलेली प्रकरणे अशी: $!A$ आण्विक आहे,
किंवा एखाद्या $\Delta_0$ !!{sentence} $!B$
साठी $\lnot !A \ident \lnot\lnot !B$.
$!A$ आण्विक असेल, तर
\olref{lem:atomic-completeness} नुसार
$\Th{Q} \Proves \lnot !A$.
$\lnot !A \ident \lnot\lnot !B$ असेल,
तर तर्कनियमांनी ते $\Th{Q}$ मध्ये
$!B$ शी सिद्धतया समतुल्य आहे.
$\lnot !A$ हे $\Struct{N}$ मध्ये सत्य
असल्याने तेही $\Struct{N}$ मध्ये सत्य आहे.
म्हणून विगमन गृहीतकाने
$\Th{Q} \Proves \lnot !A$.
\end{enumerate}
\end{proof}

\begin{prob}
\olref{lem:delta0-completeness} मधील
अस्तित्ववाची प्रकरणाचा सविस्तर पुरावा द्या.
\end{prob}

\begin{thm}
\ollabel{thm:sigma1-completeness}
$!A$ हे $\Struct{N}$ मध्ये सत्य असलेले
$\Sigma_1$ !!{sentence} असेल, तर
$\Th{Q} \Proves !A$.
\end{thm}

\begin{proof}
$\lexists[x][!A(x)]$ हे $\Struct{N}$ मध्ये
सत्य असलेले $\Sigma_1$ !!{sentence} असेल,
तर अशी नैसर्गिक संख्या~$n$ आणि चल-मूल्यन~$s$
असतात की $s(x) = n$ आणि
$\Sat{N}{!A(x)}[s]$.
पूर्ति-संबंधाच्या नेहमीच्या गुणधर्मांनुसार
$\Sat{N}{!A(\num n)}$.
परंतु $!A(\num n)$ हे $\Delta_0$ !!{formula}
आहे; म्हणून \olref{lem:delta0-completeness} ने
$\Th{Q} \Proves !A(\num n)$.
मग तर्कनियमांनी
$\Th{Q} \Proves \lexists[x][!A(x)]$.
\end{proof}

\end{document}
